% !TEX root = ../projekt.tex
% =============================================================================
% STAV: hotové podľa kódu, doplniť citácie metrík (bod 3 zadania, 1. časť).
% =============================================================================
\chapter{Metriky a metodika merania}
\label{chap:metodika}

\section{Dáta}
\label{sec:data}

Vstupom sú záznamy tokov zo siete CESNET v~CSV s~11 stĺpcami: začiatok a koniec toku,
zdrojová a cieľová IP adresa a port, protokol a počty bajtov a paketov v~oboch smeroch.
Hlavným vzorkom je \texttt{odid1}: približne 35 miliónov záznamov, ktoré po rozvinutí
dávajú približne 2,28 miliardy paketov a súčasne až približne 782 tisíc živých smerov tokov.

\TODO{Doplniť: odkiaľ dáta sú, kedy a na akej linke boli zachytené, priemerná rýchlosť v~Gb/s.}

\paragraph{Známe vlastnosti dát.}
\begin{itemize}
  \item \textbf{Zrkadlené duplicity.} \texttt{odid1} obsahuje zrkadlené záznamy tokov
        (približne 12\,\% riadkov a 42\,\% paketov). Na výreze dát nadhodnocovali preťaženie
        o~2--4\,\%, poradie funkcií nemenili. Výsledky v~tejto práci sú zatiaľ
        \textbf{s~duplicitami}. \TODO{Rozhodnúť: finálne čísla na \texttt{odid1\_dedup}, alebo to uviesť ako obmedzenie.}
  \item \textbf{Okrajový efekt.} Na začiatku záznamu chýbajú už bežiace toky a na konci
        zostávajú iba dlhé toky. Pri Toeplitze, váhe \texttt{packet} a 32 kanáloch je priemerné
        preťaženie v~prvých 10\,\% okien 138,1\,\textperthousand, v~stredných 80\,\% 90,2
        a v~posledných 10\,\% 137,1\,\textperthousand.
\end{itemize}

\section{Rekonštrukcia prúdu paketov}
\label{sec:rekonstrukcia}

Dáta popisujú toky, RSS však pracuje s~paketmi. Každý smer toku sa preto rozvinie na pakety
rozložené rovnomerne medzi začiatkom a koncom toku. Paket $k$ z~$N$ má čas
\begin{equation}
  t_k = t_\mathit{start} + (t_\mathit{end} - t_\mathit{start}) \cdot \frac{k}{N-1} .
\end{equation}
Pakety všetkých tokov sa spájajú podľa času pomocou minimovej haldy. Halda obsahuje iba
živé toky, preto pamäť nerastie s~počtom paketov.

\emph{Obmedzenie:} skutočné pakety rovnomerne rozložené nie sú, zhluky v~rámci toku
simulácia nevidí.

\section{Váha paketu}
\label{sec:vahy}

\begin{description}
  \item[flow] každý smer toku je jedna jednotka; meria kvalitu hashu bez vplyvu veľkosti toku,
  \item[packet] každý paket má váhu 1,
  \item[byte] paket váži priemernú veľkosť paketu svojho toku.
\end{description}

\section{Metriky}
\label{sec:metriky}

Nech $x_i$ je záťaž kanála~$i$ v~okne, $n$ počet kanálov a $W=\sum_i x_i$ celková záťaž okna.

\paragraph{Preťaženie} je podiel záťaže nad spravodlivým podielom, v~promile:
\begin{equation}
  O = \frac{\sum_{i} \max\left(0,\, x_i - W/n\right)}{W/1000} .
\end{equation}
Hodnota 0 znamená dokonale rovnomerné rozdelenie. Prakticky je to podiel prevádzky, ktorý by
bolo treba presunúť z~preťažených kanálov, aby boli všetky rovnako zaťažené.

\paragraph{Jainov index férovosti}~\TODO{citácia}
\begin{equation}
  J = \frac{\left(\sum_i x_i\right)^2}{n \sum_i x_i^2} , \qquad J \in \left[\tfrac{1}{n}, 1\right].
\end{equation}

\paragraph{Normalizované $\chi^2$}, kde $e = W/n$:
\begin{equation}
  \chi^2_\mathit{norm} = \frac{1}{W} \sum_i \frac{(x_i - e)^2}{e} .
\end{equation}

\section{Okná merania}
\label{sec:okna}

Záťaž sa vyhodnocuje v~nezávislých oknách s~pevným počtom paketov. Po každom okne sa
počítadlá vynulujú a neúplné posledné okno sa zahodí. Výsledok jednej konfigurácie je priemer
cez okná a cez 16 kľúčov. Okná sa porovnávali v~dvoch režimoch:

\begin{description}
  \item[pevné okno] 150\,000 paketov pre každý počet kanálov,
  \item[dynamické okno] 18\,750 paketov na kanál, teda $18\,750 \cdot n$ paketov
        (150\,000 pri 8 kanáloch, 2\,400\,000 pri 128).
\end{description}

Dôvod dynamického okna: pri pevnom okne pripadá pri 128 kanáloch na kanál iba približne
1\,170 paketov. Dynamické okno drží rovnaký počet paketov na kanál.

\paragraph{Zvažované a zamietnuté: výber časti tokov.}
Väčšia linka nesie viac prevádzky, preto sme skúšali simulovať menší počet kanálov na
pomernej časti tokov (8 kanálov = 1/16 tokov). Zamietli sme to, lebo by sa pri malom počte
kanálov nepoužila väčšina dát a výsledok by závisel od toho, ktoré veľké toky sa do výberu
dostanú.

\emph{Obmedzenie dynamického okna:} všetky počty kanálov vidia rovnakú prevádzku, takže
pri 128 kanáloch trvá okno 16-krát dlhšie a na kanál pripadá menej tokov.
\TODO{Overiť na datasete s~vyššou prevádzkou, ak ho vedúci poskytne.}
